problem:  
Let \((M^n,g,f)\) be a complete noncompact *expanding gradient Ricci soliton* with  
\[
\mathrm{Ric}_g + \nabla^2 f = \lambda g, \qquad \lambda < 0,
\]
and suppose that outside a compact set, \(M\) admits coordinates \((r,x)\in (R_0,\infty)\times X\) such that  
\[
g = dr^2 + r^2 h_r, \qquad h_r = h_0 + r^{-\mu} h_1 + o(r^{-\mu})
\]
for some \(\mu>0\) and smooth family of metrics \(h_r\) on a compact manifold \(X\).  Assume that the soliton potential satisfies \(f = \tfrac{\lambda}{4}r^2 + o(r^2)\) as \(r\to\infty\).  In classical regular cases, one shows that \(h_0\) is Einstein, but this uses relatively strong assumptions on the decay rate \(\mu\).  

**Problem:**  
Assume only that \(|\nabla_{g_C}^k(\Phi^*g - g_C)|_{g_C} = O(r^{-\delta-k})\) for some small \(\delta>0\), insufficient to perform a full asymptotic expansion of curvature.  
1. Do there exist complete noncompact expanding gradient Ricci solitons for which the asymptotic cone exists in this *weak sense* but whose link \(h_0\) fails to be Einstein even distributionally?  
2. Equivalently, does the inference “\(h_0\) is Einstein” survive under such low-decay or low-regularity asymptotic conical structures, or can genuinely non‑Einstein asymptotic cones occur at very weak rates?  
3. If not, can one quantify a critical decay threshold \(\delta_c>0\) separating rigidity (Einstein link) from possible non‑Einstein asymptotics?

Why is it a "good" problem:  
This version isolates the *asymptotic rigidity threshold* for expanding Ricci solitons, asking whether there can exist “roughly conical” expanders whose limiting cone fails to be Einstein when one relaxes the analytic hypotheses.  The question is new and unsolved: all current rigidity results depend crucially on fairly rapid asymptotic decay, and it is unknown whether weaker control allows non‑Einstein links.  

The problem probes the border between geometric regularity and analytic rigidity, connecting **asymptotic analysis** of geometric PDE with **stability phenomena** in Ricci flow and **analysis on singular spaces**.  It is simple to state but deeply involves elliptic regularity, weighted function spaces, and the nonlinear coupling between geometry and the soliton potential.  Resolving it would refine our understanding of how strictly the Ricci soliton equation dictates asymptotic conical geometry, establishing a new threshold phenomenon in geometric analysis.